\section{Alternative Views}
% A popular alternative view, especially in industry, is that there are very limited further innovations required to achieve an AI software engineer. In the past few years, we have seen emergent capabilities come from scaling \citep{wei2022emergent}, and this view may posit that the majority of the tasks and challenges we pose in this work can be resolved by scaling to larger models and training with more data.

An alternative view, particularly within the industry, suggests that limited further innovation is needed to achieve AI software engineers. 
This perspective argues that AI engineers, capable of contributing effectively to software development in the workplace, are already emerging or will be soon.  This is supported by the significant progress in recent years with better LLMs and agentic frameworks. 
The SWE-bench score, increasing from an initial 2\% to now 29\% as a state-of-the-art result in about a year, demonstrates this advancement. However, this view is perhaps too optimistic.  
While their has been considerable progress, current benchmark performances do not translate to real-world user experiences.
Indeed, SWE-Bench does not measure different aspects of software engineering and the involved complexities.
We believe that translating these performance gains to real-world use cases would continuous research and engineering in both improving AI system capabilities and the surrounding user-centric infrastructure.
% Real-world tasks often require broader scope, complex algorithms, and understanding of unclear requirements – aspects that current benchmarks do not fully capture.  The recent high score on various coding benchmark, while very impressive, still leaves many real-world software engineering abilities unmeasured.  
% Our view is that there are plenty of challenges ahead and continuous research is required. While current AI tools offer value, achieving the full potential of AI in software engineering requires addressing the challenges we have outlined. 
The alternative perspective, while not unreasonable given the recent progress, should not overshadow the need for more efforts in the direction we outline and AI for code as a field still has more to offer and overcome.


% alternative views: 
% - AI for code is a solved problem
% - we shouldn't be doing this because of economic reasons
% - AI safety: dangerous to have machines that can program themselves